% ============================================================
\chapter{Conclusions and Future Work}
\label{chap:conclusion}
% ============================================================

\section{Summary of findings}
\label{sec:summary}

This project asked whether replacing the correlation-based driver selection in
Hierarchical Sensitivity Parity with causal structure produces more robust
portfolios, and whether closing a performance-feedback loop helps further. It
built and validated a complete pipeline: causal factor discovery under a
directional prior (A1, C1), allocation in a causal sensitivity space that
preserves directional information inside a symmetric allocator (A2, C2), and a
closed feedback loop (A3, C3), evaluated on an approximate S\&P~100 over
2007--2024 at two windows, two discovery methods, a range of $K$, and a feedback
grid, with bootstrap significance throughout.

The results, stated objectively:
\begin{enumerate}[itemsep=2pt]
  \item \textbf{Causal driver selection helps modestly, at the short window only,
        and not robustly.} It beats correlation-HSP at one year (significantly for
        VARLiNGAM, $p=0.004$; directionally but not significantly for DYNOTEARS),
        but the edge disappears at two years and flips sign with $K$. The
        per-driver causal signal is statistically weak (zero FDR-significant
        drivers at calibration). The robust, window-independent finding is
        qualitative: causal selection rotates across regime-relevant drivers
        whereas correlation selection is tautological.
  \item \textbf{The most effective method is asset-only causal structure.} V0$'$,
        which uses \emph{no} exogenous drivers and converts an asset--asset causal
        graph directly into weights, is the best variant at one year (Sharpe
        $0.403$, significantly above both correlation-HSP and driver-based
        Causal-HSP). On this universe, the exogenous-driver machinery does not
        earn its place at the short window --- though, like everything else, this
        edge is window-dependent.
  \item \textbf{The closed loop adds nothing.} Across the entire $\alpha,\gamma$
        grid the feedback is inert, because utility can only re-rank within the
        causally-selected set --- and, even where it could, a monthly reward's
        signal-to-noise ratio of $\rsRewardSNR$ is too low to learn from. The C3
        hypothesis is rejected for this design.
  \item \textbf{The edge does not survive data-snooping correction.} Subjecting
        the \rsNtrials\ configurations to the Deflated Sharpe ratio, White's
        Reality Check, Hansen's SPA and the Model Confidence Set, the absolute
        Sharpe levels survive correction for non-normality and multiplicity, but
        the causal-over-correlation \emph{edge} does not clear a family-wise test
        and the correlation baseline \rsMcsVzeroIn\ retained in the $90\%$ Model
        Confidence Set. The second half of Q is answered in the negative.
\end{enumerate}

Stated as a direct answer to Q: causal driver selection does deliver a higher
risk-adjusted return than correlation selection, but only at the short window,
only modestly, and not in a way that survives correction for the multiplicity of
strategies tried --- while the strongest short-window result comes from
\emph{asset-only} causal structure with no exogenous drivers at all.

The blunt overall assessment is that, for this framework and universe, causal
methods are \emph{mildly and conditionally} effective, not robustly so: the only
statistically clean, economically modest win is asset-only causal structure at the
one-year window. This is a more useful scientific outcome than an overstated
``causal beats correlation'' headline would have been, precisely because the
robustness checks (two windows, two methods, the $K$-sweep, and the reproducibility
fix) were built to find fragility --- and found it. The project's contribution is
therefore a careful, reproducible \emph{characterisation} of where causal structure
helps portfolio construction and where it does not, together with a novel
weights-directly-from-graph result (V0$'$) that the closest prior work
\cite{howard2025causal} did not reach.

\section{Limitations}
\label{sec:limitations}

The limitations are intrinsic to the claims and are stated without hedging.
\begin{itemize}[itemsep=2pt]
  \item \textbf{Small, mostly sub-significant effects.} The headline edges are
        $0.01$--$0.03$ in Sharpe; only two cells are significant and only at one
        window, and the study is plausibly underpowered for effects this small
        (Section~\ref{sec:believe}). Multiple comparisons across the configuration
        grid further discount the significant cells.
  \item \textbf{Window- and $K$-fragility.} No positive result survives the move
        from a one-year to a two-year window, and the driver-based edge is not
        robust to $K$. Robustness across regime \emph{conditioning} is good;
        robustness across these \emph{design} choices is not.
  \item \textbf{Long-only equity, single asset class.} Every variant absorbs the
        full $\approx 51\%$ GFC drawdown; the work concerns relative
        driver-selection quality within a long-only equity allocator, not absolute
        return or drawdown control, and does not test other asset classes where
        cross-asset causal structure might be richer.
  \item \textbf{Universe approximation.} The top-by-market-cap reconstruction of
        the S\&P~100 lacks the official index's sector-balancing discretion, and a
        fixed snapshot-union universe carries mild residual survivorship effects.
  \item \textbf{Weak per-driver causal signal.} The calibration finding of zero
        FDR-significant drivers means the ``causal'' driver scores, at this
        universe size, are not far above noise --- a caveat for any
        driver-level interpretation.
  \item \textbf{Edge does not clear data-snooping correction.} The single-
        comparison significance of the short-window asset-only and VARLiNGAM
        results does not survive the family-wise battery
        (Section~\ref{sec:believe-rigour}): the baseline is retained in the Model
        Confidence Set, so the comparative claim in Q is, strictly, unproven.
\end{itemize}

\section{Future work}
\label{sec:future}

Several directions follow directly. \textbf{(1) Full non-linear discovery.} The
NTS-NOTEARS probe (Section~\ref{sec:nts}) showed the prior transfers and the
method is a genuinely different lens; a full backtest (now feasible only with more
compute or a smaller universe) could test whether non-linear causal structure
changes the picture. \textbf{(2) Foreground the asset-only route.} Given that
V0$'$ is the strongest variant, the natural follow-up is to develop asset-only
Causal-HRP in its own right --- richer asset-level graph summaries, alternative
graph-to-distance maps, and a test of whether its short-window edge generalises to
other universes. \textbf{(3) Network-density regimes.} Conditioning on the
discovered graph's density (in the manner of Howard et al.\
\cite{howard2025causal}) needs only a cheap discovery-only re-run that persists
per-window density, and would sharpen the regime analysis. \textbf{(4) Larger,
multi-asset universes}, where the per-driver signal may strengthen and where
cross-asset causal structure is plausibly more informative, addressing the
weak-signal limitation. \textbf{(5) Latent-variable and path-dependent
extensions}, to relax the causal-sufficiency assumption that no unobserved common
cause sits behind the drivers and assets. Across all of these, the reproducibility
infrastructure built here (the deterministic pipeline and the discovery cache)
makes such sweeps cheap to run honestly --- which, as this project found, is a
prerequisite for trusting the answers.
